\documentclass[11pt]{article}
\usepackage[margin=1in]{geometry}
\usepackage{amsmath,amssymb,bm}
\usepackage{booktabs}
\usepackage{enumitem}
\usepackage{xcolor}
\usepackage[colorlinks=true,linkcolor=blue!60!black,citecolor=blue!45!black]{hyperref}
\usepackage{microtype}
\usepackage{graphicx}
\graphicspath{{figs/}}

\newcommand{\E}{\mathbb{E}}
\newcommand{\sg}{\mathrm{sg}}
\newcommand{\R}{\mathbb{R}}
\newcommand{\numgood}[1]{\textbf{\textcolor{green!45!black}{#1}}}
\newcommand{\numbad}[1]{\textbf{\textcolor{red!65!black}{#1}}}
\newcommand{\whyq}[1]{\medskip\noindent\textbf{\textcolor{blue!50!black}{Q: Why #1}}\\}
\newcommand{\exppar}{\noindent\emph{Experiment.}\ }
\newcommand{\findpar}{\noindent\emph{Findings.}\ }

\title{Understanding Bilinear Value Functions in Goal-Conditioned RL:\\
a Socratic Record of Every Experiment\\
\large Living document --- v1, 2026-07-16}
\author{Factored-FB project}
\date{}

\begin{document}
\maketitle

\noindent\emph{Method of this document: each paragraph poses a question, describes the
experiment run to answer it (with training details at conference-paper level), and states the
findings. The chain continues until the only remaining questions are the currently
unanswerable ones about why bilinear $Q$-networks work --- collected in
\S\ref{sec:open} with proposed experiments. \S\ref{sec:related} surveys bilinearity and
factored value representations in RL. Companion: \texttt{fb-variants-and-grounding.tex}.}

\section{Common experimental substrate}\label{sec:substrate}

\textbf{Benchmark.} OGBench \texttt{cube-single-play-v0}: a 6-DoF arm, one cube,
28-D state observations, 5-D actions in $[-1,1]$; offline dataset of 1M transitions from a
scripted play policy (repeated pick-and-place at random locations). Evaluation: the 5
\texttt{singletask} goals, horizon 200, success = cube within 4\,cm of the target (the
environment's criterion). Unless stated: mean success over 50 episodes/task, 3--13 seeds,
500k gradient steps. Cross-domain cells use \texttt{antmaze-medium-navigate-v0} and
\texttt{scene-play-v0}.

\textbf{Value learner (GC-IQL).} Rewards $r\in\{-1,0\}$ ($0$ at the goal), masks
$m=1-\mathbf 1[s'{=}g]$; expectile value loss and MSE critic loss
\begin{equation}
\mathcal L_V = \E\,\ell^2_{\tau}\!\big(\min_i \bar Q_i(s,a,g)-V(s,g)\big),\quad
\mathcal L_Q = \sum_i\E\big(Q_i(s,a,g)-(r+\gamma m V(s',g))\big)^2,
\label{eq:iql}
\end{equation}
$\tau{=}0.9$, $\gamma{=}0.99$, twin critics, target EMA $\tau_{\mathrm{ema}}{=}0.005$
(critic only). Networks: 3-layer 512 MLPs, GELU + LayerNorm; Adam $3{\times}10^{-4}$, batch
256. Goal relabeling (value loss): 20\% current / 50\% geometric future ($p{=}1{-}\gamma$) /
30\% random; actor goals: uniform future on the same trajectory.

\textbf{The two heads under comparison.}
\begin{equation}
\text{monolithic (MLP):}\ Q_\theta([s,a,g])\qquad\text{vs.}\qquad
\text{bilinear:}\ Q(s,a,g)=\frac{\phi(s,a)^\top\psi(g)}{\sqrt d},\ d{=}50.
\label{eq:heads}
\end{equation}
Both use the same trunk widths; the bilinear head has $\approx 2\times$ the trunk parameters
(two towers) --- a confound with a queued control (\S\ref{sec:open}).

\textbf{Policy extraction.} (i) \emph{Flow} (FQL-style): a goal-agnostic BC flow $v(s,x_t,t)$
trained by flow matching; a one-shot actor $\pi(s,c,\varepsilon)$ trained by
$-\min_iQ_i/\sg|Q| + 3.0\,\|\pi-\mathrm{Euler}_{10}(v)\|^2$. (ii) \emph{DDPG+BC}: Gaussian
policy, $-Q/\sg|Q| - \alpha\log\pi_\beta$, $\alpha{=}0.3$.

\section{The Socratic chain}\label{sec:chain}

\whyq{do unsupervised/zero-shot RL methods fail on manipulation at all?}
\exppar Five methods --- FB, PSM, RLDP, CRL, and GCIQL as the supervised goal-conditioned
anchor --- each trained on the same dataset (10 seeds, 1M steps for the torch cells), plus a
\emph{phase diagnostic} that segments every evaluation rollout into approach $\to$ control
acquisition $\to$ maintained control $\to$ transport using cube--effector geometry and the
contact channel (3 seeds $\times$ 100 episodes per method).
\findpar GCIQL 0.72 $\gg$ FB 0.42--0.52 $>$ PSM 0.36 $>$ CRL 0.31 $>$ RLDP 0.17.
Approach failures are negligible for every method ($P(\mathrm{reach}){=}1.0$); grasping
differs by $\le 6$ points; \textbf{the entire method spread is post-grasp}:
$P(\mathrm{success}\mid\mathrm{secured})$ = FB 0.44 / MLP-GCIQL 0.80 / bilinear-GCIQL 0.98.
The object to explain is per-step action selection while in control, not reaching or
grasping.

\whyq{does FB fail --- architecture or training objective?}
\exppar The \emph{graft lattice}: hold the substrate of \S\ref{sec:substrate} fixed and swap
single ingredients between FB and GCIQL (3 seeds $\times$ 500k per cell, $\approx$40 cells).
Direction A grafts FB structure onto the scalar-TD grounding: the bilinear head
\eqref{eq:heads}; a $z$-conditioned forward map $\phi(s,a,\sg\psi(g))$; a two-tower $\phi$;
sphere normalization; FB's orthonormality penalty; and the full stack. Direction B grafts
scalar-TD ingredients into the measure framework (measure-only, task-normed residuals,
rank-1 test projections, $z$-mixture sampling). A separate \emph{annealing} family
warm-starts from the converged 0.904 checkpoint and schedules loss mixtures (fully described
in the companion document).
\findpar Structure survives the grounding swap: bilinear \numgood{0.904}, $z$-cond $\phi$
0.844, two-tower 0.864, full FB stack 0.871 (only ortho $\lambda{=}1000$ kills: 0.072).
The identical architectures under the measure objective: \numbad{0.08--0.34}. Decisive:
warm-start \emph{at} the 0.93 solution and fade the scalar TD out $\Rightarrow$ 0.05 ---
\textbf{the measure objective's fixed point on cube is bad; the objective is the whole gap.}
Also: fitting the measure well is insufficient --- the $z$-conditioned graft cut the measure
TD residual $24\times$ (6347$\to$266) and moved success only to 0.255.

\whyq{is FB's task-inference interface also broken?}
\exppar E1a: take the trained bilinear-GCIQL critic with a $z$-conditioned actor
($c=\sg[\psi(g)]$, group G4, $d{=}512$, 1M steps, 12 seeds at the anchor), and evaluate FB's
zero-shot rule: relabel 50k dataset transitions with the task reward, infer
$z=\sum_i r_i\psi(s'_i)$, rescale to the norm scale the actor saw in training, roll out with
$z$ only (the policy never sees $g$). Ablate the inference: moment vs whitened
(GLS) vs exact $\psi(g)$; $\sqrt d$ vs matched norm.
\findpar Moment inference \numgood{0.955} $=$ the exact-$z$ anchor. Norm hygiene is worth
0.29--0.46 (the $\sqrt d$ convention collapses two tasks); whitening is unnecessary and
ill-conditioned ($\psi$ covariance condition number $7{\times}10^{14}$). \textbf{The FB API
is intact; only the grounding was broken.} (An earlier ``inferred-$z\to0.07$'' result was an
on-policy sampling artifact: a badly conditioned policy rarely reaches the goal, so its own
rollouts give a degenerate $z$.)

\whyq{is the bilinear win an out-of-distribution generalization story?}
\exppar The \emph{OOD grid}: byte-identical start states and goals across all cells
(serialized to disk), crossing start distribution (dataset states vs \texttt{env.reset()},
which samples arm poses the play policy never produces --- median whitened NN-distance to
the data $2.4\times$ a held-out state's) with goal distribution (hindsight vs evaluation
task): C1 (ID/ID), C2 (ID/OOD-goal), C4 (OOD-start/ID), C3 (both = the real benchmark).
H$=$200, the environment's own success rule; the harness reproduces published numbers on C3
(MLP 0.747 vs published 0.717; bilinear 0.947 vs 0.904; FB 0.40 vs 0.42). 3 seeds.
\findpar Fig.~\ref{fig:oodgrid}. (a) \textbf{FB fails in-distribution}: C1 $=$ 0.43 while
grasping 100\% of episodes --- post-grasp steering, not generalization. (b) The bilinear--MLP
separation is the \emph{goal} axis: the goal shift is free for bilinear
($\Delta{\approx}{+}0.06$) and costs the MLP 16--24 points; the start shift hurts both
($-7$ to $-28$) and separates nobody. (c) The MLP is the \emph{better} critic fully
in-distribution (C1: 0.98 vs 0.87). Goal generalization is real but cannot explain (c) ---
keep asking.

\begin{figure}[h]\centering
\includegraphics[width=0.7\textwidth]{fig_oodgrid}
\caption{The OOD grid. FB's C1 uses the $z{=}B(g)$ route; its C4 was not run.}
\label{fig:oodgrid}
\end{figure}

\whyq{is the bilinear critic simply a more accurate value function?}
\exppar The \emph{counterfactual value probe}: reset the simulator to 200 in-hand grasp
states (bit-exact prefix replay --- snapshot/restore drifts because the IK-target integrator
lives outside \texttt{mjData}), roll out the agent's own flow policy to success/failure, and
measure the rank-AUC of the critic's value against the realized outcome. Actor held fixed
across critics. A \emph{decomposition probe} splits $Q_m = U_m(s,g) + A_m(s,a,g)$,
$U_m=\E_{a\sim\mu}Q_m$ over 256 behavior-flow actions, and evaluates each part separately.
\findpar MLP AUC \numgood{0.81}; bilinear \numbad{0.39}, below chance. The MLP's entire
signal is in the action-independent baseline $U$ (0.82); the bilinear's $U$ is
anti-calibrated (0.40); across states the two baselines are uncorrelated (Spearman $-0.04$)
while the two \emph{advantage} orderings partially agree ($+0.22$). \textbf{Not accuracy:
the two heads allocate capacity to different parts of $Q$} --- MLP to the baseline, bilinear
to the action-dependent part. (This also revises an earlier program thesis: cross-state
value fidelity is neither necessary nor what carries control.)

\whyq{what, then, does the bilinear critic do better?}
\exppar \emph{Reranking}: freeze the trained critic; at every environment step draw $N$
candidate actions from a frozen goal-agnostic BC flow $\mu$ and execute
$\arg\max_iQ(s,a_i,g)$. C3 protocol, $N\in\{1..128\}$, 50 episodes/task. Controls: swap the
candidate flow between agents; \emph{fresh-actor} experiments train a new flow actor (or a
new DDPG+BC actor) from scratch against the frozen critic (500k steps, optimizer masking
verified bit-frozen).
\findpar Fig.~\ref{fig:rerank}. Bilinear-as-ranker \numgood{0.91--0.96} at $N{=}32$ $=$ its
fully trained actor (0.947), \emph{with no actor training}; declines to 0.85 at $N{=}128$
(more draws find the off-support peak, next rung). MLP-as-ranker \numbad{0.57--0.59} $<$ its
trained actor (0.747). FB-as-ranker \emph{declines with $N$} (0.52 $\to$ 0.17): its $Q$
actively misleads about actions. Fresh actors reproduce the gap (flow: 0.88 vs 0.68;
co-adaptation refuted). \textbf{The bilinear critic's product is the on-support action
ordering; the ordering is the policy.} Corroboration: action-dependence is $<1\%$ of $Q$
variance for every critic (bilinear 0.46\%, FB 0.21\%, MLP 0.10\%) --- the loss barely sees
it; only the bilinear head turns it into usable structure.

\begin{figure}[h]\centering
\includegraphics[width=0.6\textwidth]{fig_rerank_nsweep}
\caption{Frozen-critic reranking of $N$ behavior-flow candidates (the mechanism figure).}
\label{fig:rerank}
\end{figure}

\whyq{does the gradient-ascent actor lose with a critic whose ordering is near-perfect?}
\exppar \emph{Off-reference test}: for the co-trained DDPG+BC agents, measure (i) the
whitened nearest-neighbor distance from the deployed actor's action to a 256-action
behavior-reference cloud, (ii) $Q(a_{\mathrm{DDPG}})-\max_{a\in\mathrm{ref}}Q(a)$ and the
fraction of states where it is positive, (iii) action-gradient norms
$\|\nabla_a Q\|$ (equal-slope vs peak discrimination). Fresh-DDPG+BC-vs-frozen-bilinear
isolates intrinsicness.
\findpar The bilinear critic values the fabricated action above \emph{every} reference
action 79\% of the time (median $+2.14$); the MLP 6\% ($-0.26$). Gradient norms are equal
(3.18 vs 3.06): an off-support \emph{maximum}, not a steeper landscape. The DDPG+BC actor
departs the behavior manifold $10\times$ further on bilinear (whitened distance 17.9 vs
1.76); a fresh DDPG+BC against the frozen critic lands at 0.52 ($\approx$ co-trained 0.557
$\ll$ flow 0.88). \textbf{The critic's one flaw lives exactly where gradient extraction
looks and anchored extraction never does} --- producing the actor reversal (bilinear+flow
0.947 / bilinear+DDPG+BC 0.51 / MLP $\approx$0.73 either way).

\whyq{is the off-support peak intrinsic to the bilinear form?}
\exppar T3: recompute the peak statistic on the bilinear+DDPG+BC checkpoints of three
domains (cube, antmaze-medium, scene; 512 dataset states $\times$ 5 eval goals; reference
cloud = the dataset action at $s$ plus the 64-NN states' actions in whitened observation
space --- no BC flow needed), plus a dataset-only action-coverage statistic
(dimension-corrected against a same-size uniform null).
\findpar Fig.~\ref{fig:peak}. Cube 0.639 (replicating the $\mu$-cloud version), antmaze
\numgood{0.003}, scene 0.014 --- the peak does not exist on antmaze, which is why
bilinear+DDPG+BC \emph{wins} there (0.837 vs MLP 0.695). The pre-registered
\emph{data-coverage} explanation is \numbad{refuted}: cube's conditional action manifold is
only $1.11\times$ narrower than antmaze's --- far too small for a $0.639\to0.003$ flip.
\textbf{The peak's location is a property of the domain's value geometry, not the data}
(open thread, \S\ref{sec:open}).

\begin{figure}[h]\centering
\includegraphics[width=0.4\textwidth]{fig_t3_peak}
\caption{The off-support peak by domain.}
\label{fig:peak}
\end{figure}

\whyq{does the bilinear head rank actions well while the MLP ranks barely above chance ---
is the MLP's action signal merely noisy?}
\exppar T1, \emph{goal-smoothed reranking}: rerank with
$\bar Q_K(s,a,g)=\frac1K\sum_{j}Q(s,a,g_j)$ where the $g_j$ perturb the goal's cube-position
dimensions by $\mathcal N(0,\sigma^2)$, $K\in\{1,8,32\}$, $\sigma\in\{2,4\}$\,cm
($=$ the success radius), 250 episodes/arm; control arm: $g_j$ = random dataset goals
(smoothing must be \emph{local} to help). Gate: $K{=}1$ must reproduce the known rerank
numbers (it does: 0.924 / 0.616).
\findpar Fig.~\ref{fig:t1}. Smoothing makes the MLP monotonically \emph{worse}
(0.616 $\to$ 0.408 at $K{=}32,\sigma{=}4$\,cm); a 2--4\,cm goal perturbation is already
distribution shift for the MLP head (the same brittleness as its $-16$ to $-24$ goal-shift
loss). Bilinear is invariant at 2\,cm (0.912) and both collapse under the random-goal
control (0.148/0.040), validating the instrument. \textbf{The MLP's action signal is absent
or systematically wrong, not noisy} --- eval-time pooling cannot recover it. If cross-goal
pooling is the bilinear mechanism, it acts at \emph{training} time, through the shared
$\phi$ basis. The causal control for that hypothesis is open (\S\ref{sec:open}).

\begin{figure}[h]\centering
\includegraphics[width=0.62\textwidth]{fig_t1_smoothing}
\caption{T1: eval-time goal smoothing does not rescue the MLP.}
\label{fig:t1}
\end{figure}

\whyq{is the factorization also responsible for training stability?}
\exppar Ablate the bottleneck: \texttt{bilinear\_sg} replaces $\psi(g)$ with $\psi(s,g)$
(same parameter count up to the input layer; every other knob fixed; 3 seeds). T2
telemetry on the resulting checkpoints and on 50k-interval retrains of the stock heads:
normalized feature co-adaptation $\E[\phi(s,a)^\top\phi(s',a')]$ (the DR3 statistic),
per-leg norms, $q$ extremes, on a fixed 4096-transition batch.
\findpar 2/3 seeds diverge late (value loss $\times$10--30; $Q$ crosses positive on a
$\le0$-bounded value); the stable seed lands at MLP level (0.69). The DR3 statistic does
\emph{not} discriminate (saturated $\approx$0.99 for all heads); the divergence lives in
the \textbf{state-conditioned goal leg}: $\E\|\psi\|^2 = 1083/2355/3789$ across seeds,
monotone in divergence severity, $q_{\max}$ to $+1.2{\times}10^3$, while $\phi$ stays sane.
\textbf{Revised mechanism: global coupling as implicit regularization} --- with $\psi(g)$,
any score inflation must pass through globally shared features and is corrected by the rest
of the data; $\psi(s,g)$ opens a local $(s,g)$ pocket that expectile-TD's upward pressure
(no output bound, no $V$ target) inflates.

\whyq{how much rank does the win need, and which incidental knobs matter?}
\exppar Rank sweep $d\in\{5,28,50,128,512\}$ and self-ablations (single critic head;
expectile $\tau\in\{0.5,0.9,0.95\}$), 3--13 seeds each.
\findpar $d{=}5$ collapses (0.12); $d\ge 28$ is flat (0.86--0.91): \textbf{rank is a floor,
not a dial} --- the task family's demands are low-rank; the win is not about tuning $d$.
Single-head 0.874 ($-3$); $\tau{=}0.5$ \numgood{0.920}, $\tau{=}0.95$ 0.874 --- the expectile
is not the ingredient either ($\tau{=}0.5$ is plain SARSA-style expectile and works
slightly better).

\whyq{does the effect transfer to other value learners and domains?}
\exppar The SOTA replication: swap MLP$\to$bilinear inside GCIVL (AWR extraction), HIQL
(hierarchical), and across antmaze/scene; 3--8 seeds per cell.
\findpar Cube: GCIVL 0.675 vs 0.496 stock; HIQL improves only when the bilinear TD loss is
routed \emph{around} the shared subgoal representation (two-head split 0.35--0.42 vs shared
0.15--0.19; the published-tuning cell is confounded and was rerun at the correct
\texttt{subgoal\_steps}). Antmaze: bilinear+DDPG+BC \numgood{0.837} (vs 0.695); flow fails
there for \emph{both} critics (0.24--0.33) --- an actor-side artifact (the goal-agnostic BC
anchor fights direction commitment on a direction-multimodal marginal), currently being
fixed by conditioning the anchor. Scene: bilinear loses with both actors (0.26--0.29 vs
0.52) and exhibits \emph{no} off-support peak --- a distinct, undiagnosed failure
(\S\ref{sec:open}). \textbf{The critic-side win transfers across value learners on cube;
the extraction layer is what breaks portability.}

\whyq{summing up --- why does bilinear GCIQL + flow work on cube?}
Three verified clauses. (1) \emph{Sound grounding}: scalar hindsight TD is contractive and
task-weighted; the measure alternative is neither, and its fixed point here is bad.
(2) \emph{Capacity reallocation}: the factorization cannot ignore the $<1\%$-variance action
signal the way a monolithic net can --- one shared $\phi$ serving every goal accumulates it
into a near-perfect on-support ordering, paying with value calibration (irrelevant
downstream), gaining free goal generalization and backup stability, and leaving one
structured flaw (the off-support peak). (3) \emph{Extraction match}: anchored extraction
consumes exactly the ordering and never visits the flaw; gradient extraction does the
opposite. Cube is where all three trades pay simultaneously (goal-shift-dominated
benchmark; post-grasp ordering bottleneck; benign anchor).

\section{The bedrock: open questions about bilinear $Q$-networks}\label{sec:open}

These are the questions the chain could not answer; each carries its proposed experiment.
This section is the live frontier of the document.

\begin{enumerate}[leftmargin=1.6em,itemsep=3pt]
\item \textbf{Why does the low-rank constraint force a good action ordering?} The surviving
hypothesis is training-time cross-goal pooling through the shared $\phi$ basis; the rivals
(multimodality, direction-vs-distance content, eval-time denoising) are dead. \emph{Queued
(T5)}: (a) a \emph{late-fusion} control --- identical two towers, MLP readout on
$[\phi;\psi]$ instead of the inner product (isolates the readout from the tower structure);
(b) a parameter-matched MLP; (c) feature-rank (srank) trajectories of the MLP critic to test
the rank-collapse reading of its near-chance ordering.
\item \textbf{Where does the off-support peak form?} Coverage is refuted; the working idea
is quasimetric-residual relocation (a symmetric rank-$d$ form cannot represent the
asymmetric optimal goal-value; the residual must land somewhere, and its location follows
the domain's value geometry). \emph{Queued}: fit the peak's location against embedding-norm
growth off-manifold ($\|\phi(s,a)\|$ along $a_{\mathrm{ref}}\to a_{\mathrm{DDPG}}$ rays),
cube vs antmaze; a positive result yields a \emph{predictive} account (and a test:
norm-capped $\phi$ should shrink the peak).
\item \textbf{Why does the same head lose on scene?} No peak, both actors fail; hypotheses:
combinatorial goal structure under a single rank-50 $\psi$; or the HIQL-style shared-rep
poisoning. \emph{Queued}: scene rerank N-sweep (in flight), then a Schaul-style split and a
per-factor goal head.
\item \textbf{How low-rank is the \emph{true} advantage?} $d\ge28$ sufficing suggests the
optimal action-ordering structure of cube is genuinely low-rank; is it? \emph{Queued (T9)}:
estimate $\mathrm{rank}(A^*(s,\cdot,g))$ over the goal grid from counterfactual rollouts
(atlas banks + simulator).
\item \textbf{Regime inversion vs the online literature.} Bilinear Value Networks (online,
HER, DDPG) found the \emph{state-conditioned} pairing $f(s,a)^\top\phi(s,g)$ essential;
offline we find the opposite ($\psi(g)$ stable, $\psi(s,g)$ divergent). Is the difference
the offline expectile backup (no bound + no $V$-target + local inflation pockets), or the
data regime? \emph{Queued}: $\psi(s,g)$ with an output bound / $V$-target --- if it
stabilizes at bilinear level, the divergence, not the pairing, is the offline story.
\item \textbf{Does the ordering quality have a clean information-theoretic signature?}
The critic's $U$ and $A$ decouple (uncorrelated baselines, partially shared advantages).
\emph{Speculative}: measure mutual information between $\phi$'s top principal subspace and
one-step controllable state deltas, bilinear vs MLP vs FB --- a representation-level
correlate of the reranking gap that could be computed at any checkpoint.
\end{enumerate}

\section{Related work: bilinearity and factored value representations in RL}
\label{sec:related}

\emph{(Living section; entries carry one-line relevance to our findings.)}

\paragraph{Successor structure.} Successor representations (Dayan 1993) and successor
features $\psi^\pi(s,a)$ with $Q=\psi^\top w$ (Barreto et al.\ 2017, arXiv 1606.05312)
are the original ``value = features $\times$ task-vector'' factorizations; USFA (Borsa et
al.\ 2019, 1812.07626) adds policy-conditioning and GPI. FB (Touati \& Ollivier 2021;
Touati, Rapin, Ollivier 2023, 2209.14935) factorizes the successor \emph{measure}
$M\approx F^\top B\rho$ --- our failed grounding. Our bilinear GCIQL is, in this taxonomy,
successor features whose reward basis $\psi(g)$ is co-learned by scalar hindsight TD ---
``SFs without feature TD'' (every explicit vector-TD estimator of the same structure
collapsed in our lattice: 0.00--0.38).

\paragraph{Bilinear value heads.} Bilinear Value Networks (Hong, Yang, Agrawal, ICLR 2022,
2204.13695) is the closest prior: $Q=f(s,a)^\top\phi(s,g)$ inside online DDPG+HER, with
gains attributed to compositional goal generalization --- matching our free-goal-shift
finding, but with the pairing inverted relative to our stable configuration
(\S\ref{sec:open}, item 5). FLAP (2101.04750) shows shared-basis linear heads extrapolate
across tasks in meta-RL. Low-rank structure of $Q$ matrices: Yang et al.\ (ICLR 2020,
1909.12255) and Shah, Song, Yu (NeurIPS 2020) verify empirical control $Q$-matrices are
near-low-rank and use rank restriction as a structured denoiser --- independent support for
``rank as a helpful prior, not a capacity loss'' ($d\ge28$ flat in our sweep).

\paragraph{Embedding-factored values from representation learning.} Contrastive RL
(Eysenbach et al.\ 2022, 2206.07568) learns $f(s,a,g)=\phi(s,a)^\top\psi(g)$ as a
goal-conditioned critic by InfoNCE --- structurally our head, trained by a different
(contrastive, on-support-ranking-shaped) objective; C-learning (2011.08909) is its
TD-flavored sibling. VIP (Ma et al., ICLR 2023, 2210.00030) derives an \emph{implicit}
time-contrastive objective from the Fenchel dual of goal-conditioned RL and produces a value
as embedding distance $V(s;g)=-\|\varphi(s)-\varphi(g)\|$, pre-trained from passive video
--- evidence that value-shaped embedding geometry can be learned without TD at all; LIV
(2306.00958) extends it to language--image goals. HILP (Park et al., ICML 2024, 2402.15567)
embeds temporal distances isometrically ($\|\varphi(s)-\varphi(g)\|=d^*(s,g)$); METRA
(Park, Rybkin, Levine, ICLR 2024, 2310.08887) maximizes $(\varphi(s')-\varphi(s))^\top z$
under a temporal-metric constraint --- both are inner-product/distance task-factored values
with explicit geometric constraints where ours emerges from TD alone. Laplacian
representations (Wu, Tucker, Nachum, ICLR 2019) and spectral decompositions of transition
operators (Ren et al.\ 2022, 2208.09515) motivate the same low-rank structure from operator
theory.

\paragraph{Quasimetric obstructions.} $-V^*$ of goal-reaching is a quasimetric; symmetric
inner products cannot represent it with low distortion --- QRL (Wang, Torralba, Isola,
Zhang, ICML 2023, 2304.01203), IQE (2211.15120), MRN (Liu, Zhu, Stone, AAAI 2023,
2208.08133). This is the principled source of our ``irreducible, structured residual'' ---
and the open question of \emph{where} it lands (\S\ref{sec:open}, item 2).

\paragraph{Theory: linear and low-rank MDPs.} Linear MDPs (Jin, Yang, Wang, Jordan, COLT
2020, 1907.05388) make $Q^\pi$ exactly linear in known features $\phi(s,a)$ for every
policy; FLAMBE (Agarwal et al.\ 2020, 2006.10814) and representation learning for low-rank
MDPs (Uehara, Zhang, Sun, ICLR 2022, 2110.04652) learn the factorization. Our setting is
the ``representation-learning'' regime of this theory with $\psi(g)$ as a learned task
embedding rather than a known reward parameter.

\paragraph{TD and representation dynamics.} DR3 (Kumar et al., ICLR 2022, 2112.04716):
TD implicitly grows feature dot-products of bootstrap-adjacent pairs (our T2
operationalization of it saturated and did not discriminate; the divergence we measure is a
norm/pocket phenomenon instead). Implicit under-parameterization (Kumar et al., ICLR 2021,
2010.14498): bootstrapping collapses feature rank --- the standing candidate for the MLP's
absent action signal (srank probe queued). Self-predictive learning dynamics (Tang et al.,
ICML 2023, 2212.03319) analyze when two-tower TD-style objectives collapse --- relevant to
why our SF vector-TD cells die while scalar TD through the same structure lives.
Norm/scale-constrained critics: HL-Gauss/Stop Regressing (Farebrother et al., ICML 2024,
2403.03950) and SimbaV2 (2502.15280) --- bounded/structured critic outputs beating raw MSE
regression under TD; our $\sqrt d$-scaled inner product (and the T10 scale-anchor episode)
are instances of the same principle.

\paragraph{Extraction.} One-step RL (Brandfonbrener et al., NeurIPS 2021), IDQL
(2304.10573), SfBC (2209.14548), FQL (Park, Li, Levine 2025, 2502.02538), and the
one-step$\equiv$critic-regularization equivalence (Eysenbach et al.\ 2023, 2307.12968)
together predict our extraction findings: selection/anchored consumption of a frozen
well-ordering critic is near-optimal offline; gradient consumption amplifies critic
artifacts (BCQ, 1812.02900).

\paragraph{Benchmark.} OGBench (Park, Frans, Eysenbach, Levine, ICLR 2025, 2410.20092).
Published FB-on-cube anchors: TD-JEPA (2510.00739) reports FB 62.0 (RGB) on cube-single and
$\le$3 on cube-double; no published work diagnoses \emph{why} --- this document's chain is,
to our knowledge, the only mechanism-level account.

\section{Experiment queue (living)}\label{sec:queue}
T5 late-fusion + param-matched controls (open item 1) $\cdot$ peak-geometry norm probe
(item 2) $\cdot$ scene rerank + split heads (item 3) $\cdot$ T9 true-advantage rank
(item 4) $\cdot$ bounded-$\psi(s,g)$ (item 5) $\cdot$ srank trajectories (item 1c) $\cdot$
multi-env rerank N-sweeps (antmaze/scene; partially in flight) $\cdot$ RG-FB d512 ladder
run $\cdot$ T10 500k reads.

\end{document}
